\chapter{Introduction}
\label{ch:intro}

\begin{guidance}[What this chapter does]
Sets the project in context and leads the reader to clearly stated aims and
objectives. By the end of this chapter the marker should know
\emph{what} problem you tackled, \emph{why} it matters (technically,
commercially, for users or society) and \emph{how} the report is organised.
Typical length: 600--900 words.

Write the report in the past tense for the work you carried out (\enquote{The
prototype was tested \ldots}) and in the present tense for what the report
contains (\enquote{\Cref{ch:results} presents \ldots}).
\end{guidance}

\section{Background}
\label{sec:background}

\begin{guidance}
Introduce the subject area and the problem. Support claims with references,
for example \textcite{engc2020ahep} or \parencite{cross2008}. Include
commercial considerations where relevant: who would use or buy the
solution, and what it would replace.
\end{guidance}

Optimisation of energy use in residential homes is an active area of
research \parencite{bausch2014}. \todo{Replace with your own background.}

\section{Problem statement}
\label{sec:problem}

\begin{guidance}
One or two paragraphs that end in a precise statement of the problem. Some
projects frame this as a single research question, for example:
\begin{quote}
How can complicated real-world applications be built with an easy way of
programming microcontrollers?
\end{quote}
(from an engineering diploma thesis, University of Portsmouth, 2008).
\end{guidance}

\section{Aim and objectives}
\label{sec:objectives}

\begin{guidance}
\textbf{One aim}: the overall outcome of the project.
\textbf{Several objectives}: specific, measurable steps that together achieve
the aim. Each objective should start with a verb and be checkable at the end
of the project. You will come back to every objective in
\cref{ch:conclusions}, so number them with the \texttt{objectives} list and
refer to them with \verb|\cref{obj:...}|.

Example taken from a PhD thesis \parencite{bausch2014}:
\emph{Aim}: to create a prediction model for inside air temperatures in a
vacant residential home that uses environmental and experimental data from
disparate sources. \emph{Objectives}: to create a system that collects data
from local, remote and aggregated sensors; to create prediction models from
historical data; \ldots
\end{guidance}

The aim of this project was to \ldots

The objectives were:
\begin{objectives}
  \item\label{obj:review} To review existing solutions and the relevant
    literature to establish the requirements for \ldots
  \item\label{obj:spec} To produce a product design specification \ldots
  \item\label{obj:design} To generate and evaluate candidate solutions and
    select the most suitable one using a structured method.
  \item\label{obj:build} To implement and test a prototype \ldots
  \item\label{obj:evaluate} To evaluate the prototype against the
    specification and the project process against the plan.
\end{objectives}

\section{Scope and constraints}
\label{sec:scope}

\begin{guidance}
What is in, what is out, and why. Constraints include budget, time,
equipment, safety, standards and access to data.
\end{guidance}

\section{Structure of the report}
\label{sec:structure}

\begin{guidance}
One or two sentences per chapter. Use \verb|\cref| so that the numbers update
automatically.
\end{guidance}

\Cref{ch:literature} reviews the literature and existing products.
\Cref{ch:methodology} describes the project methodology and plan, and
\cref{ch:specification} the product design specification. \Cref{ch:design}
develops and evaluates candidate solutions. \Cref{ch:implementation}
describes the implementation and testing, \cref{ch:issues} the issues that
arose, and \cref{ch:results} the results. \Cref{ch:evaluation} evaluates the
product and the process, \cref{ch:conclusions} draws conclusions against the
objectives, and \cref{ch:recommendations} makes recommendations.
